\documentclass[12pt]{article}
\usepackage[margin=1in]{geometry}
\usepackage{ctex}
\usepackage{graphicx}
\usepackage{caption}
\usepackage{subcaption}
\usepackage{hyperref}
\usepackage{listings}
\usepackage{xcolor}
\usepackage{fontspec}  % 如果需要使用等宽字体
\usepackage{float}

% 定义颜色
\definecolor{codegreen}{rgb}{0,0.6,0}
\definecolor{codegray}{rgb}{0.5,0.5,0.5}
\definecolor{codepurple}{rgb}{0.58,0,0.82}
\definecolor{backcolour}{rgb}{0.95,0.95,0.92}

% 现代简约风格
\lstdefinestyle{mystyle}{
    backgroundcolor=\color{backcolour},   
    commentstyle=\color{codegreen},
    keywordstyle=\color{magenta},
    numberstyle=\tiny\color{codegray},
    stringstyle=\color{codepurple},
    basicstyle=\ttfamily\footnotesize,
    breakatwhitespace=false,         
    breaklines=true,                 
    captionpos=b,                    
    keepspaces=true,                 
    numbers=left,                    
    numbersep=5pt,                  
    showspaces=false,                
    showstringspaces=false,
    showtabs=false,                  
    tabsize=2,
    frame=single,
    framerule=0.5pt,
    rulecolor=\color{gray!30},
    xleftmargin=15pt,
    xrightmargin=5pt
}

\lstset{style=mystyle}
\hypersetup{colorlinks=true, linkcolor=blue, urlcolor=blue, citecolor=blue}

\title{雀宝 MahjongPaw 项目文档}
\author{栾仲文 \& 顾宸尧}
\date{\today}

\graphicspath{{img/}}

\begin{document}
\maketitle
\tableofcontents
\newpage

\section{项目概述}
本项目旨在为日式立直麻将提供桌面端 AI 助手能力，集成图像识别、牌理分析与大模型流式解说，支持一键截图到策略建议的完整闭环。

\section{项目整体架构}
总体流程包括输入截图、YOLO 感知、麻将内核分析、LLM 解说与 GUI 展示五部分，形成“截图→识别→分析→解说”的闭环。

\begin{lstlisting}
MahjongPaw/
├── main.py                    # 项目入口（启动 GUI）
├── api_gui.py                 # GUI 主界面（流式输出、聊天、Markdown 渲染）
├── analyzer.py                # 牌局分析器（向听数、听牌、打点计算）
├── api.py                     # 大模型 API 调用（DeepSeek）
├── cal_scores.py              # 点数计算模块
├── Mahjong_YOLO/              # YOLO 麻将牌识别模块
│   ├── trained_models_v2/     # 训练好的 YOLO 模型权重
│   │   ├── yolo11n_best.pt
│   │   ├── yolo11s_best.pt
│   │   └── yolo11m_best.pt
│   ├── notebooks/             # 数据标注和处理
├── world_model/               # 世界模型（麻将对象定义）
│   ├── mahjong_table.py       # 麻将桌类
│   ├── mahjong_player.py      # 玩家类
│   ├── mahjong_tile.py        # 麻将牌类
│   └── mahjong_meld.py        # 副露类
├── models/                    # 模型文件目录
│   └── yolov8_mahjong.pt      # YOLO 模型权重（备用）
├── requirements.txt           # 项目依赖
└── README.md                 # 项目说明文档
\end{lstlisting}

\section{项目模块详述}
\subsection{日式立直麻将内核}
负责向听数计算、听牌张检索、打点评估与役种统计，作为策略分析的基础。

\begin{lstlisting}[caption={麻将内核demo}]
class MahjongTable:
    """"
    游戏桌类：
    管理场次类型（半庄战）、
    局次类型（南一局）、
    庄家信息（用户1）、
    手牌区与副露区、
    牌河区、
    牌山区、
    岭上区、
    宝牌区、
    玩家点数（25000点）、
    立直信息、
    校验全局信息（全局点数和守恒）
    """
    pass

    
class Player:
    """"
    玩家类：
    玩家信息、
    当前点数、
    是否庄家、
    是否立直、
    手牌与副露、
    切/吃/碰/杠/胡方法
    """
    pass
    
    
class Hand:
    """"
    手牌类：
    暗手子类、
    副露子类、
    存储、排序、查找、移除等方法
    """
    pass
    
    
class MahjongTile:
    """
    麻将牌类，用于表示一张麻将牌的属性和行为
    包含类型：万、条、筒（序数牌）；东、南、西、北（风牌）；中、发、白（箭牌）
    """
    pass
\end{lstlisting}

\subsection{YOLO图片识别}
负责从游戏截图中识别玩家手牌并序列化为标准牌谱字符串，为后续牌理计算提供输入。

\begin{figure}[H]
    \centering
    \includegraphics[width=\textwidth]{test.png}
    \caption{测试截图}
\end{figure}

\begin{figure}[H]
    \centering
    \includegraphics[width=\textwidth]{prediction_result.jpg}
    \caption{模型输出示例}
\end{figure}

\subsection{牌效计算}
结合听牌张与役种，评估翻符与点数，排序输出高打点方案与建议。

模块耦合度较高，分布在不同的文件中，共同作用计算出有效进张的牌效，其结果被保存到output.txt中。

\begin{lstlisting}[caption={牌效计算模块（部分）}]
from mahjong.tile import TilesConverter
from mahjong.hand_calculating.hand import HandCalculator
from mahjong.hand_calculating.hand_config import HandConfig
from mahjong.meld import Meld
from mahjong.shanten import Shanten
from mahjong.agari import Agari
from cal_scores import calc_hand_score
from Mahjong_YOLO.test import perceive
from collections import Counter
import numpy as np
\end{lstlisting}

\begin{lstlisting}[caption={牌效计算模块输出示例}]
======================================================================
🀄 麻将手牌分析报告
======================================================================

📊 基础信息:
   手牌: 566m777p1244668s
   牌数: 13张
   自风: 东
   场风: 东
   立直: 否

🎯 牌型分布:
    万子: 3张
    筒子: 3张
    索子: 7张
    字牌: 0张
    总牌数: 13张
    幺九牌: 1张 (7.7%)

🎲 向听数分析:
    普通和牌向听数: 2
    🔄 2向听

🎯 听牌分析:
    ❌ 未听牌

======================================================================
🎉 分析完成！
======================================================================
\end{lstlisting}

\subsection{大模型API调用，及GUI界面}
通过大模型进行流式解说，GUI 承载交互与展示，一键完成截图、识别、分析与解说。

\begin{figure}[H]
    \centering
    \includegraphics[width=\textwidth]{structure.png}
    \caption{桌面 GUI 主界面示意}
\end{figure}

\section{项目主要界面贴图或者主要实验结果}
\subsection{GUI 主界面}

\begin{figure}[H]
    \centering
    \includegraphics[width=\textwidth]{structure.png}
    \caption{桌面 GUI 主界面示意}
\end{figure}

GUI使用Tkinter库实现，主界面包括：
\begin{enumerate}
    \item 一键截图+分析+AI解说：一键截图，获取雀宝评论。
    \item 雀宝聊天框：雀宝的输出文本会实时显示在聊天框中。
    \item 聊天窗口：对于雀宝的分析，哪里没听懂的还可以继续追问。
    \item 开始分析（流式输出）：开发者玩具。
    \item 清空结果：这个按键可以清空结果。  
\end{enumerate}

\subsection{麻将牌图像识别效果对比}

根据上传截图，当前玩家手牌为\textbf{68p256777889s34z}。

麻将专业大模型雀宝的输出为\textbf{68p256777889s34z}，与上传截图中的手牌一致；
但是业内卓越的多模态大模型豆包的输出为\textbf{222m256p25688s34z}，答案错误。

事实上，将手牌全部识别正确是一项极其困难的任务。对于13张手牌的识别，需要有90\%以上的准确率才能被认为是准确的,
这意味着单张牌面的识别准确度要超过99.2\%。

\begin{figure}[H]
    \centering
    \includegraphics[width=\textwidth]{doubao_kan.png}
    \caption{多模态大模型豆包的输出}
\end{figure}

\begin{figure}[H]
    \centering
    \includegraphics[width=\textwidth]{real_kan.png}
    \caption{麻将专业大模型雀宝的输出}
\end{figure}

注：以上对比基于手牌\textbf{68p256777889s34z}，其他手牌的识别结果可能会有差异。
雀宝使用了YOLOv8n模型，该模型在训练时使用了10000张麻将牌图片，每张图片包含13张牌面，
模型的输入为3通道的224x224像素图片，输出为13个类别的概率分布。而豆包（我认为）使用的是Transformer模型（在多模态方面有十分好的效果）。
但是在准确率和召回率方面，不如配合监督学习的YOLOv8n模型。

\subsection{麻将向听计算效果对比}

手牌\textbf{566m777p1244668s}向听数应该为2，即需要至少两张有效进张才能听牌。

而麻将专业大模型雀宝的输出为2，答案正确。一般的但是业内领先的非推理大模型豆包的输出为1，答案错误。

这说明麻将专业大模型雀宝在向听计算方面的效果要优于一般的非推理大模型豆包。术业有专攻，
向听数计算是麻将专业大模型雀宝的一项核心功能，而一般的非推理大模型豆包在这方面的效果一般。

\begin{figure}[H]
    \centering
    \includegraphics[width=\textwidth]{doubao_tin.png}
    \caption{多模态大模型豆包的输出}
\end{figure}

\begin{figure}[H]
    \centering
    \includegraphics[width=\textwidth]{real_tin.png}
    \caption{麻将专业大模型雀宝的输出}
\end{figure}

注：以上对比基于手牌\textbf{566m777p1244668s}，其他手牌的向听数计算结果可能会有差异。
雀宝调用的是 DeepSeek 提供的 API，而网页版DeepSeek的向听数计算准确性也劣于雀宝。
这都是基于项目Workflow中的牌效评估模块，而牌效评估模块是麻将专业大模型雀宝的一项核心功能。

\section{总结}
项目整合视觉感知、牌理计算与 LLM 解说，形成面向实战的桌面助手。后续可扩展场况特征、危险度评估与复盘功能。

\end{document}
